\PassOptionsToPackage{main=ngerman}{babel}
%\PassOptionsToPackage{numbers}{natbib}
%\PassOptionsToPackage{main=english}{babel}

\documentclass[thesis]{mi-document}

% Microtypografie aktivieren
\usepackage{microtype}


\bachelor % im Falle einer Masterarbeit \master

% Variablen, die für das Deckblatt und Metadaten verwendet werden
\title{Evaluierung von Einsatzmöglichkeiten einer Deep Learning Applikation zur Detektion gelabelter Merkmale in Bildern}
\author{Amina Kasa}
\semester{Sommersemester 2024}
\course{Bachelorarbeit}
\module{Fakultät Informatik und Mathematik}
\dozent{Prof. Dr. Timo Baumann}
\studid{3190755}
\studSemester{Medizinische Informatik}
\phone{0941/133742666} % Optional
\studSubject{Medizinische Informatik}
\firstReviewer{Prof. Dr. Timo Baumann}
\secondReviewer{[Zweitgutachter]}
\advisor{GEFASOFT Automatisierung und Software GmbH}
\address{Klenzestr. 2a, 93051 Regensburg}{} % Optional
\mail{kasaamina@gmail.com}
\studMail{amina.kasa@st.oth-regensburg.de}
\dateHandedIn{[Abgabetermin der Arbeit]}
\keywords{Enter;key;words;here}

%\usepackage[numbers]{natbib}
\bibliographystyle{apacite}
% Falls Sie die Abkürzung zum Einbinden von Grafiken benutzen möchten. Erläuterung fnden Sie im Abschnitt zu Abbildungen.
\newcommand*{\image}[2]{
	\begin{figure}
	\centering
	\includegraphics[width=0.5\textwidth]{{images/#1}}
	\caption{#2}
	\label{fig:#1}
	\end{figure}
}

\usepackage{emptypage}
\pagenumbering{roman}

\usepackage[acronym]{glossaries}
\makeglossaries

\begin{document}
% ---------------------------------------------------------------
\pagestyle{empty}
% Das Deckblatt erstellen
 

% \newlength\headpicwd
% \setlength{\headpicwd}{0.2\pdfpagewidth} % image width of physical page?
% \newcommand\printlogo{%
%   \AddToShipoutPictureBG*{%
%     \AtStockUpperLeft{%
%       \raisebox{-2cm}{%
%       \hspace{2.5cm}
%         % centred with respect to physical page ..?
%         % \hspace*{\dimexpr 0.5\pdfpagewidth - 0.5\headpicwd\relax}%
%         \includegraphics[keepaspectratio=true, width=0.33\pdfpagewidth]{images/template/unili_logo.jpg}%
%       }%
%     }%
%   }%
% }

\begin{titlepage}
\setstretch{1.5}

\vspace*{-30pt}
\includegraphics[width=0.5\textwidth]{images/Unbenannt}
\hfill\includegraphics[width=0.5\textwidth]{images/logo_firma}
\vspace{20pt}

\centering
\Large
\setstretch{1.5}
\textbf{\getInstitute}\\
\large
\textbf{\getModule}\\
\normalsize
\textbf{Studiengang \getStudSubject}

\vspace{60pt}
\setstretch{1.5}
\huge
\textbf{Exposé zur}

%\vspace{5pt}
\setstretch{1.5}
\Huge
\textbf{\getWorkType}

\normalsize

\textmd{\getAuthor}
\vspace{20pt}

\Large
\textbf{\getTitle}

\vspace{20pt}

\normalsize


%\ifcurrentbaselanguage{German}{
%\getWorkType im \getStudSubject{}\\
%at \getInstitute
%}

%\ifcurrentbaselanguage{English}{
%\getWorkType in \getStudSubject{}\\
%at \getInstitute
%}

\setstretch{1.5}

\vfill

\begin{tabularx}{\linewidth}{@{}l<{:}>{\RaggedRight\arraybackslash}X@{}}
\ifcurrentbaselanguage{German}{

Vorgelegt von&\getAuthor\\
Matrikelnummer&\getStudID\\
E-Mail &\getStudMail\\
Erstgutachter&\getFirstReviewer\\
Zweitgutachter&\getSecondReviewer\\
Betreuende Firma&\getAdvisor\\
Abgegeben am&\getDateHandedIn

}

\end{tabularx}
\end{titlepage}


% \thispagestyle{empty}
% \quad \addtocounter{page}{-1}
% \newpage
% \pagenumbering{roman}   


 \newpage

 % Abstract hinzufügen
%    \clearpage
%    \doublespacing

%    \begin{abstract}[ngerman]

\end{abstract}
%    \clearpage 

% ---------------------------------------------------------------
%Inhaltsverzeichnisse  
\pagestyle{plain}
\singlespacing
\addcontentsline{toc}{chapter}{Inhaltsverzeichnis}
\tableofcontents % Optional
\newpage

% Abkürzungsverzeichnis
\clearpage
\addcontentsline{toc}{chapter}{Abkürzungsverzeichnis}
\printglossary[type=\acronymtype, title=Abkürzungsverzeichnis]

% \addcontentsline{toc}{chapter}{Abbildungsverzeichnis}
% \listoffigures % Optional
% \newpage

% \addcontentsline{toc}{chapter}{Tabellenverzeichnis}
% \listoftables % Optional
% \newpage

% \addcontentsline{toc}{chapter}{Quellcodeverzeichnis}
% \lstlistoflistings % Optional
% \newpage
% --------------------------------------------------------------

\newacronym{ai}{AI}{Artificial Intelligence}
\newacronym{ibv}{IBV}{industrielle Bildverarbeitung}
\newacronym{dl}{DL}{Deep Learning}
\newacronym{cnn}{CNNs}{Convolutional Neural Networks}
\newacronym{knn}{KNNs}{Künstlichen Neuronalen Netzen}
\newacronym{bv}{BV}{Bildverarbeitung}
\newacronym{cv}{CV}{Computer Vision}
\newacronym{oe}{OE}{Objekterkennung}
\newacronym{ml}{ML}{Maschinelles Lernen}

\pagenumbering{arabic}
\pagestyle{headings} % Seitennummern und Kapitelbezeichnungen anzeigen
    
    % hier beginnt der eigentliche Inhalt der Arbeit
    \chapter{Einleitung}\label{sec:exp_einleitung}
Die industrielle Bildverarbeitung hat sich zu einer Schlüsseltechnologie in modernen Fertigungsprozessen entwickelt. Sie ermöglicht die automatisierte Inspektion, Qualitätskontrolle und Prozessoptimierung, wodurch Effizienzsteigerungen und Kostensenkungen erreicht werden können. Durch den Einsatz von Kamerasystemen und spezialisierter Software werden visuelle Informationen erfasst und analysiert, um beispielsweise Defekte frühzeitig zu erkennen, die Position von Objekten zu bestimmen oder Produktionsabläufe zu überwachen (Cognex, 2023). (https://www.cognex.com/de-de/what-is/machine-vision)

Mit der zunehmenden Komplexität industrieller Produkte und Prozesse steigen jedoch auch die Anforderungen an Bildverarbeitungssysteme. Besonders herausfordernd sind Szenarien, in denen Bilder durch Rauschen oder andere Störeffekte beeinträchtigt sind. Eigenschaften wie spezielle Texturen, Formen, Farben oder Größen von Objekten können die zuverlässige Detektion erschweren, da Fehler mit dem Hintergrund verschmelzen und somit eine präzise Analyse verhindern. Bildrauschen kann die Genauigkeit der Objektdetektion erheblich beeinträchtigen, was zu fehlerhaften Inspektionen und potenziellen Qualitätsproblemen führt.

Convolutional Neural Networks (CNNs) haben sich als leistungsfähige Werkzeuge in der Bildverarbeitung etabliert, da sie komplexe Muster und Strukturen in Bilddaten erkennen können \cite{lecun_deep_2015}. Ein besonderer Fokus liegt auf der semantischen Segmentierung, bei der jedem Pixel eines Bildes eine Klasse zugewiesen wird, was eine detaillierte Analyse und Interpretation ermöglicht \cite{long_fully_2015}.

Das U-Net-Modell ist eine spezialisierte Architektur von CNNs, die ursprünglich für die biomedizinische Bildsegmentierung entwickelt wurde \cite{ronneberger_u-net_2015}. Aufgrund seiner Fähigkeit, auch mit begrenzten Daten und in schwierigen Umgebungen präzise Ergebnisse zu liefern, hat es sich in verschiedenen Anwendungsbereichen bewährt. Die Anwendung von U-Net in industriellen Kontexten, insbesondere unter herausfordernden Bedingungen wie starkem Bildrauschen, bietet vielversprechende Möglichkeiten, die Leistungsfähigkeit von Bildverarbeitungssystemen zu steigern.

\section{Motivation}
Während meines Praxissemesters bei der Gefasoft GmbH in Regensburg hatte ich die Gelegenheit, an Projekten der industriellen Bildverarbeitung zu arbeiten. Dabei kamen verschiedene traditionelle Methoden zum Einsatz, wie Schwellenwertverfahren, Kantendetektion, Regionenwachstum, Watershed-Algorithmus und Clustering-Verfahren. Diese Methoden sind unerlässlich für eine erfolgreiche Bildsegmentierung, Objektdetektion und Bildklassifikation. \cite{suse_bildverarbeitung_2014,szeliski_computer_2022}Die Grundidee besteht darin, Bildinformationen zu vereinfachen, um nachfolgende Analysen zu erleichtern und die Komplexität zu reduzieren.Erfolgreiche Segmentierung, die für eine effiziente Bildanalyse entscheidend ist, beinhaltet die Aufteilung des Bildes in homogene Regionen. Sie spielt eine wesentliche Rolle bei der Objekterkennung und der Abgrenzung von Grenzen, indem Pixeln mit gemeinsamen visuellen Merkmalen identifiziert werden können.\cite{jahne_digitale_2024}.Diese Ansätze sind attraktiv, da sie keine komplexen Architekturen oder hohen Rechenaufwand erfordern und oft leicht zu implementieren sind. 

Jedoch wurde ich bei der Anwendung dieser Methoden auch mit deren Herausforderungen bekannt. Insbesondere bei stark verrauschten Bildern und geringem Kontrast zwischen Defekt und Hintergrund lieferten sie keine zufriedenstellenden Ergebnisse. Feine Defekte, wurden häufig nicht erkannt, da sie sich kaum vom Hintergrund unterschieden und das Rauschen die relevanten Signale überlagerte.Die traditionellen Methoden sind stark auf manuell gestaltete Merkmale und Heuristiken angewiesen, was arbeitsintensiv ist und ihre Anpassungsfähigkeit an neue oder veränderte Bedingungen einschränkt. Zudem erreichen sie im Vergleich zu modernen Verfahren oft eine geringere Genauigkeit und Effizienz, insbesondere bei komplexen Aufgabenstellungen. \cite{zhao_object_2019} -> Anmerkung vielleicht ein paar Bilder und Fälle die dir begegnet sind.

Diese Einschränkungen sind für meine Arbeit von besonderer Bedeutung, da die zuverlässige Erkennung von Defekten in stark verrauschten Bildern für die Qualitätssicherung in industriellen Prozessen entscheidend ist.

Im Gegensatz zu traditionellen Methoden lernen Deep-Learning-methoden, ins besondere Convolutional Neuronal Networks (CNNs),  automatisch, komplexe Merkmale aus Bildern zu extrahieren. Diese Fähigkeit ermöglicht es den Modellen, sich ohne manuelle Merkmalsauswahl an eine Vielzahl von Aufgaben anzupassen, was zu robusteren und skalierbareren Lösungen führt \cite{mo_review_2022}. Deep-Learning-Modelle bieten in der Regel eine höhere Genauigkeit bei Aufgaben wie Objektdetektion, Segmentierung und Klassifikation, da sie aus großen Datensätzen lernen können. Sie können komplexe Muster in den Daten erfassen, die von traditionellen Methoden oft übersehen werden \cite{zhao_object_2019}. Zudem können Deep-Learning-Modelle so trainiert werden, dass sie robust gegenüber Variationen in Bildqualität, Beleuchtung und Hintergrundrauschen sind, was häufige Probleme bei traditionellen Bildverarbeitungsmethoden sind \cite{liu_deep_2023}.Diese Vorteile sind direkt relevant für die Herausforderungen, denen ich in meiner praktischen Arbeit begegnet bin.

\section{Problemstellung}
Das Training neuronaler Netze erfordert üblicherweise eine große Anzahl annotierter Bilder mit geeignetem Kontext. In der industriellen Qualitätskontrolle steht jedoch oft nur eine begrenzte Anzahl spezifischer, fehlerhafter und fehlerfreier Musterteile zur Verfügung. Da es für jeden Anwendungsfall keine öffentlich zugänglichen Datensätze gibt, entsteht ein relativ kleiner Datensatz, dessen zu detektierende Regionen manuell annotiert werden müssen. Dieser Mangel an umfangreichen und vielfältigen Trainingsdaten stellt eine erhebliche Herausforderung dar.

Diese Einschränkungen führen dazu, dass für das Training eines Deep-Learning-Modells nur sehr wenige Bilder verfügbar sind, was Probleme wie Overfitting begünstigen kann. Das Modell neigt dazu, die Trainingsdaten zu überanpassen und generalisiert schlecht auf neue, unbekannte Daten.\cite{lecun_deep_2015}

Vor diesem Hintergrund stellt sich die Frage, wie eine effektive semantische Segmentierung in der industriellen Qualitätskontrolle trotz begrenzter Datenmenge erreicht werden kann.

Eine vielversprechende Lösung bietet die U-Net-Architektur, die sich in der Literatur als geeignet für die semantische Segmentierung komplexer Bilddaten erwiesen hat. U-Net stützt sich intensiv auf Datenaugmentation, um die begrenzte Anzahl annotierter Beispiele effizient zu nutzen. Das Netzwerk wird mit sehr wenigen Bildern trainiert und zeigt dennoch eine überlegene Leistung bei Aufgaben der biomedizinischen Bildsegmentierung \cite{ronneberger_u-net_2015}. Diese Eigenschaften machen U-Net zu einem idealen Kandidaten für Anwendungen in der industriellen Qualitätskontrolle, bei denen nur eingeschränkte Trainingsdaten verfügbar sind.

In den letzten Jahren wurde die U-Net-Architektur erfolgreich in verschiedenen Bereichen der Medizin eingesetzt\cite{azad_medical_2024,siddique_u-net_2021}. Ihre Anpassungsfähigkeit und Effizienz bei der Verarbeitung komplexer Bilddaten unter schwierigen Bedingungen legen nahe, dass sie auch für die Herausforderungen in der industriellen Bildverarbeitung geeignet sein könnte.

    %\chapter{Elevator Pitch}\label{sec:exp_elevator}
In meiner Bachelorarbeit untersuche ich den Einsatz von \gls{dl}-Methoden, speziell der U-Net-Architektur, zur automatischen Detektion gelabelter Merkmale in Bildern. Traditionelle Bildverarbeitungsmethoden stoßen bei stark verrauschten Bildern und begrenzten Trainingsdaten oft an ihre Grenzen. Durch die Implementierung eines U-Net-Modells möchte ich zeigen, wie trotz dieser Herausforderungen eine präzise und effiziente Erkennung relevanter Merkmale erreicht werden kann. Mein Ziel ist es, die Möglichkeiten von Deep Learning in der industriellen Qualitätskontrolle zu evaluieren und somit zu einer verbesserten Fehlererkennung und Prozessoptimierung in Fertigungsprozessen beizutragen.








    \chapter{Forschungsfragen und Zielsetzung}\label{sec:exp_ziel}
\section{Forschungsfragen}
\begin{enumerate}
    \item Wie effektiv kann das U-Net-Modell gelabelte Merkmale in stark verrauschten industriellen Bildern mittels semantischer Segmentierung detektieren, und welche Genauigkeit kann dabei erreicht werden?
    \item Welche Techniken der Datenvorbereitung (z. B. Rauschunterdrückung, Datenaugmentation) und Modellanpassung (z. B. Architekturmodifikationen, Hyperparameter-Tuning) können die Leistungsfähigkeit des U-Net-Modells unter Rauscheinflüssen verbessern?
\end{enumerate}

\section{Zielsetzung}
Das Ziel dieser Bachelorarbeit ist es, die Segmentierung von Merkmalen in Bildern mithilfe einer U-Net-Architektur zu automatisieren. Es werden segmentierte Bilder verwendent, um Defekte oder relevante Objekte zu erkennen und ihre Grenzen präzise zu bestimmen. Die U-Net-Architektur ist ein universeller Funktionsapproximator mit hoher Genauigkeit und eignet sich besonders gut für verrauschte Eingabedaten (Chen et al., 2016). Ursprünglich für die biomedizinische Bildsegmentierung entwickelt, könnte U-Net auch in industriellen Anwendungen von großem Nutzen sein, insbesondere beim Umgang mit komplexen Bilddaten. Das Ziel meiner Arbeit ist es daher, diese Methode zu untersuchen, um die Analyse dieser Bilddatensätze zu verbessern und gleichzeitig zeitaufwendige manuelle Prozesse zu minimieren oder zu vermeiden.


Das Ziel dieser Bachelorarbeit ist es, die U-Net-Architektur für die semantische Segmentierung in industriellen Bildern zu untersuchen und zu optimieren. Ein zentraler Bestandteil der Arbeit ist die Einarbeitung in den Stand der Technik durch eine umfassende Literaturrecherche, um aktuelle Trends und Herausforderungen im Bereich der Deep-Learning-basierten Bildverarbeitung zu verstehen.

Darüber hinaus wird ein besonderer Fokus auf die Präparierung der Bilddaten gelegt, indem Bilddatensätze sorgfältig annotiert und durch Techniken wie Datenaugmentation erweitert werden. Hierbei kommt verfügbare Bildverarbeitungssoftware zum Einsatz, um die Qualität der Trainingsdaten zu maximieren.

Ein weiterer Schritt ist die Auswahl geeigneter Softwarebibliotheken, wobei die Implementierung des Modells in PyTorch vorgenommen wird. Die optimale Modellbildung und Anpassung an spezifische Anforderungen der industriellen Bildverarbeitung stehen im Vordergrund.

Abschließend erfolgt die Validierung des Modells durch eine systematische Bewertung der Ergebnisse. Die Arbeit zielt darauf ab, die Effizienz und Genauigkeit des U-Net-Modells zu bewerten und es mit traditionellen Methoden der Bildverarbeitung zu vergleichen, um die besten Strategien für die industrielle Qualitätskontrolle zu identifizieren.

    \chapter{Methodik}\label{sec:exp_methodik}
\cite{arafin_deep_2024} -->> sehr wichtige Quelle zur realisierung der Arbeit.
\begin{enumerate}
    \item Datenbeschaffung und -vorbereitung:
    
    \begin{itemize}
        \item Datensätze: Verwendung von industriellen Bilddatensätzen mit relevanten Objekten und Defekten aus dem Betrieb. Es werden konkrete Beispiele verwendet, die in der Praxis mittels traditioneller Segmentierungs methoden, keine gewünschten ergebnisse erziehlt haben.
        \item Labeling: Die Bereiche der Bilder die den Fehler zeigen, werden für jeden Bild händisch aufgezeichnet. Dies erfordet hintergrund wissen, über die gesuchten Fehlern in den Bildern, da sie auch bei blosem Auge sehr schwer zu erkennen sind.
        \item Datenvorbereitung: Es werden Methoden wie Binning und aufteilung in Patches angewendet.
    \end{itemize}
    
    \item Modellimplementierung:
    \begin{itemize}
        \item Software: Die Implementierung erfolgt in Python unter Verwendung von Bibliotheken wie PyTorch für das Modelltraining sowie Matplotlib zur Visualisierung von Ergebnissen.
        \item Modellanpassung: Die U-Net-Architektur wird durch gezielte Modifikationen der Layer-Struktur angepasst, um sowohl die Effizienz als auch die Geschwindigkeit des Trainings zu verbessern.
        \item Hyperparameter-Tuning: Eine systematische Optimierung von Hyperparametern wie Lernrate, Batch-Größe und weiteren Modellparametern wird durchgeführt, um die Leistung des Modells zu maximieren.
    \end{itemize}
    
    \item Training und Validierung:
    \begin{itemize}
        \item Trainingsstrategie: Ein UNet-Modell wird implementiert, das auf der Kombination verschiedener Verlustfunktionen, wie Binary Cross Entropy und Dice Loss, basiert. Die Verwendung von Lernraten-Schedulern, wie dem StepLR, spielt eine zentrale Rolle, um die Effizienz und den Fortschritt des Trainingsprozesses zu optimieren.
        \item Datenaugmentation: Techniken wie Rotation, Spiegelung, Skalierung und Helligkeitsvariation zur Datenaugmentation angewendet, um das Modell robuster zu machen und die Generalisierung zu verbessern 
        \item Regularisierung: Regularisierungstechniken wie Dropout oder Early Stopping integriert, um das Modell besser gegen Überanpassung (Overfitting) abzusichern und die allgemeine Performance auf unbekannten Daten zu verbessern.
    \end{itemize}

    \item Evaluation:
    \begin{itemize}
        \item Metriken: Bewertung anhand von Metriken wie Intersection over Union (IoU), Dice-Koeffizient, Präzision, Recall und F1-Score.
        \item Analyse der Rauschresistenz: Untersuchung der Modellleistung bei unterschiedlichen Rauschpegeln und -arten.
    \end{itemize}
    
    \item Analyse und Diskussion:
    \begin{itemize}
        \item Ergebnisinterpretation: Identifikation von Erfolgsfaktoren und Limitierungen.
        \item Praktische Implikationen: Bewertung der Anwendbarkeit in industriellen Szenarien.
        \item Empfehlungen: Entwicklung von Leitlinien für die Implementierung.
    \end{itemize}
    
\end{enumerate}
    \chapter{Vorläufige Gliederung}
\label{sec:exp_gliederung}

% Provide an expected Table of Content of the final thesis.

\preliminaryTOCbox{
% You can remove the "box" by deleting the \preliminaryTOCbox command
\begin{enumerate}[label=\arabic*,font=\bfseries] 
    \item[] \hspace{-1cm} \textbf{Abstract}\enspace  \secdotfill \enspace i
    %\item \textbf{Introduction} \secdotfill 1
    \item \textbf{Einleitung} \secdotfill  i
        \begin{enumerate}[label*=.\arabic*]
            \item Hintergrund und Motivation \subsecdotfill i
            \item Zielsetzung und Fragestellung \subsecdotfill i
            \item Aufbau der Arbeit \subsecdotfill i
        \end{enumerate}
        
    \item \textbf{Literaturrecherche und Stand der Technik} \secdotfill  i
        \begin{enumerate}[label*=.\arabic*]
            \item Überblick über die Industrielle Bildverarbeitung \subsecdotfill i
            \item Deep Learning in der Bildverarbeitung \subsecdotfill i
            \item Aktuelle Forschung und Entwicklungen \subsecdotfill i

        \end{enumerate}
    
    \item \textbf{Präparierung der Bilddaten} \secdotfill  i
        \begin{enumerate}[label*=.\arabic*]
            \item Beschreibung des Datensatzes \subsecdotfill i
            \item Labeling-Prozess \subsecdotfill i
            \item Datenaugumentationstechniken \subsecdotfill i

        \end{enumerate}

    \item \textbf{Modellierung und Training} \secdotfill  i
        \begin{enumerate}[label*=.\arabic*]
            \item Auswahl der Softwarebibliotheken \subsecdotfill i
            \item Modellarchitektur \subsecdotfill i
            \item Trainingsprozess und Hyperparameter-Optimierung \subsecdotfill i

        \end{enumerate}

    \item \textbf{Validierung und Bewertng der Ergebnisse} \secdotfill i
        \begin{enumerate}[label*=.\arabic*]
            \item  Methoden der Validierung \subsecdotfill i
            \item Bewertungsmethoden und Metriken \subsecdotfill i
            \item Analyse und Diskussion der Ergebnisse \subsecdotfill i

        \end{enumerate}

    
    \item \textbf{Fazit und Ausblick} \secdotfill  i
        \begin{enumerate}[label*=.\arabic*]
            \item  Zusammenfassung der Ergebnisse \subsecdotfill i
            \item Diskussion der Ergebnisse \subsecdotfill i
            \item Ausblick auf Zukünftige Arbeiten \subsecdotfill i

        \end{enumerate}
    
    \preliminaryTOCend
 
\end{enumerate}
}
    \chapter{Zeitlicher Ablauf}\label{sec:exp_ablauf}
\begin{table}[h!]
    \centering
    \begin{tabular}{|p{3cm}|p{3cm}|p{8cm}|}
        \hline
        \textbf{Phase} & \textbf{Zeitraum} & \textbf{Beschreibung} \\
        \hline
        \textbf{Woche 1 – 2} & Literaturrecherche & 
        - Einarbeitung in den Stand der Technik \newline
        - Lesen und Zusammenfassen relevanter Studien zu U-Net und Bildverarbeitung \\
        \hline
        \textbf{Woche 3 – 4} & Datenpräparierung & 
        - Sammeln und Annotieren von Bilddaten \newline
        - Durchführung von Datenaugmentation mit Bildverarbeitungssoftware \\
        \hline
        \textbf{Woche 5 – 6} & Modellimplementierung & 
        - Auswahl geeigneter Softwarebibliotheken, wie \textit{PyTorch} \newline
        - Implementierung der U-Net-Architektur \\
        \hline
        \textbf{Woche 7 – 8} & Modelltraining und -anpassung & 
        - Training des U-Net-Modells auf den vorbereiteten Datensätzen \newline
        - Optimierung durch Hyperparameter-Tuning und Modellanpassung \\
        \hline
        \textbf{Woche 9-10} & Validierung & 
        - Evaluierung der Modellleistung durch Metriken wie Genauigkeit, Präzision und Recall \\
        \hline
        \textbf{Woche 11} & Ergebnisanalyse & 
        - Analyse und Interpretation der Ergebnisse \newline
        - Diskussion der Stärken und Schwächen der Methode \\
        \hline
        \textbf{Woche 12} & Fertigstellung und Abgabe & 
        - Abschluss der Bachelorarbeit \newline
        - Einholen von Feedback \newline
        - Letzte Überarbeitungen und Abgabe der Arbeit \\
        \hline
    \end{tabular}
    \caption{Zeitlicher Ablauf der Bachelorarbeit}
\end{table}
    
    % Kapitelbezeichnung in der rechten oberen Ecke entfernen
\clearpage
\pagestyle{plain}

% --------------------------------------------------------------   
% Literaturverzeichnis anzeigen
% kleinerer Zeilenabstand, damit es nicht so gestreckt aussieht
\onehalfspacing
\bibliography{references}
\doublespacing
\end{document}
